\subsection{深度与内射维数}

命题 8.--- 设 A 为 Noether 环，M 为有限生成 A-模。我们有 \(\dim_{A}(M) \leqslant \mathrm{di}_{A}(M)\)。

设 \(r\) 为一个小于 \(\dim_{\mathrm{A}}(\mathrm{M})\) 的正整数。存在素理想的一条饱和链 \(\mathfrak{p} \subset \mathfrak{p}_{1} \subset \ldots \subset \mathfrak{p}_{r-1} \subset \mathfrak{q}\)，使得 \(\mathfrak{p}\) 为 M 的支撑集的极小元；于是 \(\mathrm{A}_{\mathfrak{p}}\) -模 \(\mathrm{M}_{\mathfrak{p}}\) 具有有限长度，故有 \(\mathrm{Hom}_{\mathrm{A}_{\mathfrak{p}}}(\kappa(\mathfrak{p}), \mathrm{M}_{\mathfrak{p}}) \neq 0\)，从而 \(\mathrm{Ext}_{A_{\mathfrak{q}}}^{r}(\kappa(\mathfrak{q}), \mathrm{M}_{\mathfrak{q}}) \neq 0\) ( \(\S 1, n^{\circ} 7\) 引理 3)，这蕴含 \(\mathrm{di}_{\mathrm{A}}(\mathrm{M}) \geqslant r\) ( \(n^{\circ} 3\) , 命题 6)。

命题 9.--- 设 A 为 Noether 局部环，M 为具有有限内射维数的非零有限生成 A-模。则有 \(\mathrm{di}_{\mathrm{A}}(\mathrm{M}) = \mathrm{prof}(\mathrm{A})\)。

令 \(r = \mathrm{di}_{A}(\mathrm{M})\)。我们有 \(\operatorname{Ext}_{A}^{i}(\kappa_{A}, \mathrm{M}) = 0\) 对 \(i > r\) 成立，故有 \(\operatorname{Ext}_{\mathrm{A}}^{r}(\kappa_{\mathrm{A}}, \mathrm{M}) \neq 0\)（由 \(\mathfrak{n}^{\circ}\) 3 的命题 6, (iv)⇒(i)）。设 \(s\) 为 A 的深度，\((x_1, \ldots, x_s)\) 为 \(\mathfrak{m}_{\mathrm{A}}\) 中元素的一个 A-正则序列 (§ 1, \(\mathfrak{n}^{\circ}\) 4, 定理 2)；令 \(\mathrm{N} = \mathrm{A}/(x_1\mathrm{A} + \ldots + x_s\mathrm{A})\)。由 \(\mathfrak{n}^{\circ}\) 1 的例，我们有 \(\mathrm{dp}_{\mathrm{A}}(\mathrm{N}) = s\) 且 \(\operatorname{Ext}_{\mathrm{A}}^{s}(\mathrm{N}, \mathrm{M}) \neq 0\)，故 \(s \leqslant \mathrm{di}_{\mathrm{A}}(\mathrm{M}) = r\)。但 N 的深度为 0 (§ 1, \(\mathfrak{n}^{\circ}\) 4, 命题 7)，故存在 A-模的正合列

\[
0 \rightarrow \kappa_ {\mathrm{A}} \longrightarrow \mathrm{N} \longrightarrow \mathrm{N} ^ {\prime} \rightarrow 0.
\]

由此导出扩张模的一个正合列

\[
\operatorname{Ext} _ {\mathrm{A}} ^ {r} (\mathrm{N}, \mathrm{M}) \longrightarrow \operatorname{Ext} _ {\mathrm{A}} ^ {r} (\kappa_ {A}, \mathrm{M}) \longrightarrow \operatorname{Ext} _ {\mathrm{A}} ^ {r + 1} (\mathrm{N} ^ {\prime}, \mathrm{M});
\]

由于 \(\mathrm{Ext}_{\mathrm{A}}^{r + 1}(\mathrm{N}^{\prime},\mathrm{M}) = 0\) 且 \(\mathrm{Ext}_{\mathrm{A}}^{r}(\mathbf{k}_{\mathrm{A}},\mathrm{M})\neq 0\)，我们得到 \(\mathrm{Ext}_{\mathrm{A}}^{r}(\mathrm{N},\mathrm{M})\neq 0\)，故 \(r\leqslant \mathrm{dp}_{\mathrm{A}}(\mathrm{N}) = s\)。最终有 \(r = s\)，证明完毕。
% semantic-fragments: legacy-input-seam
\relax{}
